% Задание на УИР — оформление по образцу thesis-template-3-pz.pdf.
% Размеры шрифтов сняты непосредственно с PDF примера (см. диалог):
% "Задание на УИР" -- 17.2пт полужирный, "Студенту..." -- 14.3пт обычный,
% "ТЕМА УИР"/"ЗАДАНИЕ" -- 14.3пт полужирный, название темы -- 17.2пт
% обычный, таблица -- 10пт (колонка "Форма отчётности" -- 8пт, как в
% примере), подпись "Дата, подпись" -- 8пт.
% Содержание таблицы отражает реальную структуру разделов 1-5 данной ПЗ
% (см. main.tex): каждый пункт задания соответствует конкретному
% подразделу, фактически написанному в разделах razdel1.tex--razdel5.tex.
% Не нумеруется (управляется через \pagestyle{empty} в main.tex).

\pzletterhead

\vspace{0.5cm}
\begin{center}
{\fontsize{17.2}{20}\selectfont\bfseries Задание на УИР}

\vspace{0.35cm}
{\fontsize{14.3}{17}\selectfont Студенту гр. Б22-544 Портнову Л.М.}

\vspace{0.5cm}
{\fontsize{14.3}{17}\selectfont\bfseries ТЕМА УИР}

\vspace{0.35cm}
{\fontsize{17.2}{20}\selectfont
Разработка прототипа системы интеграции непрерывного профилирования\\
в CI/CD-пайплайн для раннего обнаружения регрессий производительности\\
на языке Go}

\vspace{0.5cm}
{\fontsize{14.3}{17}\selectfont\bfseries ЗАДАНИЕ}
\end{center}

\vspace{0.3cm}
{\fontsize{10}{12}\selectfont
\begin{longtable}{|p{1.0cm}|p{7.0cm}|>{\fontsize{8}{9.5}\selectfont}p{3.3cm}|p{1.7cm}|p{2.5cm}|}
\hline
\textbf{№ п/п} & \textbf{Содержание работы} & \textbf{Форма отчётности} & \textbf{Срок исполнения} & \textbf{Отметка о выполнении} \\
 & & & & {\fontsize{8}{9.5}\selectfont Дата, подпись} \\
\hline
\endfirsthead
\hline
\textbf{№ п/п} & \textbf{Содержание работы} & \textbf{Форма отчётности} & \textbf{Срок исполнения} & \textbf{Отметка о выполнении} \\
 & & & & {\fontsize{8}{9.5}\selectfont Дата, подпись} \\
\hline
\endhead
\hline
\endfoot
\textbf{1.} & \textbf{Аналитическая часть} & & & \\
\hline
1.1. & Обзор инструментов сбора и анализа данных о производительности в экосистеме Go & Подраздел ПЗ, список литературы & & \\
\hline
1.2. & Сравнительный анализ подходов к интеграции профилирования и обнаружения регрессий в CI/CD & Подраздел ПЗ, сравнительная таблица подходов & & \\
\hline
\textit{1.3.} & \textit{Оформление расширенного содержания пояснительной записки (РСПЗ)} & \textit{Текст РСПЗ} & & \\
\hline
\textbf{2.} & \textbf{Теоретическая часть} & & & \\
\hline
2.1. & Разработка модели процесса профилирования в CI/CD-пайплайне & Подраздел ПЗ, схема процесса & & \\
\hline
2.2. & Формализация статистического критерия обнаружения регрессии производительности & Подраздел ПЗ, формулы критерия & & \\
\hline
2.3. & Разработка концептуальной схемы интеграции двухуровневого профилирования в пайплайн & Подраздел ПЗ, схема артефактов & & \\
\hline
\textbf{3.} & \textbf{Инженерно-технологическая часть} & & & \\
\hline
3.1. & Обоснование выбора модели жизненного цикла и инструментальных средств & Подраздел ПЗ & & \\
\hline
3.2. & Проектирование архитектуры прототипа & Подраздел ПЗ, схема архитектуры & & \\
\hline
3.3. & Описание состава и структуры основных компонентов системы & Подраздел ПЗ & & \\
\hline
\textbf{4.} & \textbf{Практическая часть} & & & \\
\hline
4.1. & Разработка тестового Go-приложения с бенчмарками & Исходный код & & \\
\hline
4.2. & Реализация критерия сравнения с эталоном (\texttt{benchstat}) & Исходный код & & \\
\hline
4.3. & Реализация диагностического механизма сравнения pprof-профилей & Исходный код & & \\
\hline
4.4. & Настройка CI/CD-пайплайна с профилированием (GitHub Actions) & Конфигурация CI, исходный код & & \\
\hline
\textbf{5.} & \textbf{Экспериментальная апробация} & & & \\
\hline
5.1. & Подготовка сценариев внесения регрессий производительности & Исходный код (ветки \texttt{regression/*}) & & \\
\hline
5.2. & Проверка корректности обнаружения регрессий на управляемых сценариях & Журналы экспериментов, таблицы результатов & & \\
\hline
5.3. & Оценка накладных расходов профилирования в CI/CD & Журналы экспериментов, таблицы результатов & & \\
\hline
5.4. & Сравнение и анализ результатов апробации & Сравнительный анализ, таблицы результатов & & \\
\hline
\textit{6.} & \textit{Оформление пояснительной записки (ПЗ) и иллюстративного материала для доклада} & \textit{Текст ПЗ, презентация} & & \\
\hline
\end{longtable}
}

\vspace{1.0cm}
\noindent
\begin{tabular}{@{}p{6cm}p{2.6cm}p{3cm}p{4cm}@{}}
Дата выдачи задания: & Руководитель & \hrulefill & \underline{Климов В.В.} \\
\hrulefill & Студент & \hrulefill & \underline{Портнов Л.М.} \\
\end{tabular}
